\centerline{\Huge Условия и решения задач.}
\centerline{\bf \Huge 4 класс}
 
\problem{4.1} Расставьте скобки так, чтобы получилось верное равенство: 
$2 \cdot 74 - 2 : 2 + 16 = 8$

\hspace{3mm}

Решение:\\ 
$2 \cdot (74 - 2) : (2 + 16) = 8$ \\
Дополнительных объяснений не требуется.

\hspace{3mm}

Критерии:\\
Правильная расстановка скобок: \textbf{7 баллов.}\\
Неправильная расстановка скобок: \textbf{0 баллов.}

\hspace{3mm}

\problem{4.2} В 9:00 на остановке «Университет» стояли 4 белых автобуса. Через 10 минут между каждым автобусом припарковались по 3 синих автомобиля. Ещё через 15 минут между каждым транспортным средством встали по 2 красных трактора. Сколько стояло транспортных средств на остановке «Университет» в 9:20, если известно, что никакое транспортное средство не уезжало с остановки?

\hspace{3mm}

Ответ: 13 транспортных средств.

Решение:\\
Количество промежутков между четырьмя белыми автобусами равно трём. Значит, в 9:10 в каждый промежуток встали по 3 синих автомобиля, и всего транспортных средств на данный момент времени стало: $4 + 3 \cdot 3 = 4 + 9 = 13.$ В следующий раз изменение количества транспортных средств произошло еще через 15 минут, то есть в 9:25. Следовательно, в 9:20 их было столько же, сколько и в 9:10. Значит, в 9:20 было 13 транспортных средств на остановке «Университет».

\hspace{3mm}

Критерии:\\
Полностью верное решение со всеми обоснованиями: \textbf{7 баллов}.\\
Полностью верное решение со всеми обоснованиями, но с арифметическими ошибками: \textbf{5 баллов}.\\
Ответом является верное количество транспортных средств в 9:25: \textbf{4 балла}.\\
\newpage
\hspace{10mm}

\noindentИдеи, продвигающие ход решения, но недостаточные для получения верного ответа: \textbf{2 балла}.\\
\noindentНеверное решение или идеи, не продвигающие ход решения: \textbf{0 баллов}.\\
Верный или неверный ответ без обоснования: \textbf{0 баллов}.

\hspace{3mm}

\problem{4.3} В каждой из двадцати школ города Элисты обучаются поровну учеников. Если в каждой школе будет обучаться на 50 учеников меньше, то во всех школах останется столько учеников, сколько их раньше было суммарно в девятнадцати школах. Сколько учеников изначально обучалось в каждой школе? 

\hspace{3mm}

Ответ: 1000 учеников.

Решение:\par
\textit{I вариант}: \\
Всего во всех школах стало обучаться на 20 · 50 = 1000 учеников меньше. До этого было 20 школ с одинаковым количеством учеников, а стало 19 школ с одинаковым количеством учеников. Значит, общее количество учеников уменьшилось на одну школу и в то же время на 1000 учеников. Следовательно, в каждой школе обучается по 1000 учеников.

\textit{II вариант}: Пусть в каждой школе обучается по X учеников. Тогда составим уравнение:\\$20\cdot X - 20\cdot50 = 19\cdot X\\
X = 20\cdot50 = 1000$. \\Следовательно, в каждой школе обучается по 1000 учеников.

\hspace{3mm}

Критерии:\\
Любое верное решение не подбором и с верным обоснованием: \textbf{7 баллов.}\\
В целом верное решение с верным обоснованием, но с допущением арифметических ошибок: \textbf{5 баллов.}\\
Решение подбором с обоснованием подбора: \textbf{4 балла.}\\
Идеи, продвигающие ход решения, но недостаточные для получения верного ответа: \textbf{2 балла.}\\
Неверное решение или идеи, не продвигающие ход решения: \textbf{0 баллов.} \\
Верный или неверный ответ без обоснования: \textbf{0 баллов.}

\newpage

\hspace{10mm}

\begin{wrapfigure}{R}{0.3\textwidth}
\vspace{-20mm}
\centering
\includegraphics[width=0.35\textwidth]{Городская Олимпиада Элиста 2023/Условия/picture_4.4.png}
\caption{Рис. 1}
\end{wrapfigure}

\problem{4.4} Однажды Пюрвя вырезал по линиям клетки из клетчатого листа бумаги 4 одинаковые фигурки и решил их склеить друг с другом. У него получилась замечательная фигура (рисунок справа). Определите, как разрезать замечательную фигуру на изначальные 4 одинаковые фигурки Пюрви?

\hspace{3mm}

\begin{wrapfigure}{R}{0.3\textwidth}
\vspace{-20mm}
\centering
\includegraphics[width=0.35\textwidth]{Городская Олимпиада Элиста 2023/Решения/picture_4.4_sol.jpg}
\caption{Рис. 1}
\end{wrapfigure}

Решение:\\
Одно из разрезаний: рисунок справа снизу\\
Дополнительных объяснений не требуется.

Критерии:\\
Любое верное разрезание, возможно отличающееся от предложенного в решении: \textbf{7 баллов.}\\
Получена площадь фигурки, равная 5 клеткам: \textbf{3 балла.}\\
Получена площадь замечательной фигуры, равная 20 клеткам: \textbf{1 балл.}\\
Любое неверное разрезание (не по линиям, не из четырех фигур, не из одинаковых фигур и т.д.): \textbf{0 баллов.}

\hspace{3mm}

\problem{4.5} На шахматной доске необычного размера 23 клетки в длину и 23 клетки в ширину, состоящей из 264 белых полей и 265 чёрных, (расположенных в шахматном порядке) стоит по одной белой пешке на каждом белом поле, а на чёрных полях ничего не стоит. В какой-то момент времени Рома решил одновременно взять каждую белую пешку и переставить в какое-то соседнее по стороне поле. Какое минимальное количество чёрных полей будет не занято, если на одном поле может стоять более одной пешки?

\hspace{3mm}

Ответ: 1 черное поле.

Решение:\\
Заметим, что соседние с белой клеткой по стороне только чёрные клетки. Следовательно, Рома переставляет каждую белую пешку с белого поля только на черное. Всего белых клеток 264, значит, и белых пешек на них стоит 264, а чёрных полей 265. Следовательно, по обратному принципу Дирихле: клеток, то есть чёрных полей, больше, чем кроликов, то есть белых пешек.

\newpage

\hspace{3mm}

\noindentТаким образом, найдется хотя бы одна клетка, то есть одно чёрное поле, которое будет пустым. А значит, минимальное количество черных полей, которые будут не заняты, равно единице. (Можно доказывать без обратного принципа Дирихле через метод от противного)

\hspace{3mm}

Критерии:\\
Полностью верное решение со всеми обоснованиями: \textbf{7 баллов.}\\
Идеи, продвигающие ход решения, но недостаточные для получения верного ответа: \textbf{2 балла.}\\
Неверное решение или идеи, не продвигающие ход решения: \textbf{0 баллов.}\\
Верный или неверный ответ без обоснования: \textbf{0 баллов.}

